\section{Цель работы}

Создать динамические библиотеки и сравнить два способа их подключения:
при линковке программы и при явной загрузке во время её работы.

\section{Условие}

Вариант 16 (функции №3 и №4). Для натуральных \texttt{A} и \texttt{B}
функция \texttt{PrimeCount(A, B)} считает простые числа на отрезке
\texttt{[A, B]}, а \texttt{GCF(A, B)} вычисляет НОД.
Команды \texttt{1 A B} и \texttt{2 A B} вызывают соответствующие функции.
Команда \texttt{0} переключает реализацию в \texttt{program2}.

\section{Метод решения}

Обе библиотеки реализуют интерфейс из \texttt{lab5.h}.

\textbf{Библиотека \texttt{impl1}.}
В \texttt{PrimeCount} каждое $n \geq 2$ на отрезке проверяется на делимость
на числа от $2$ до $n-1$. Если делителей нет, счётчик увеличивается.
В \texttt{GCF} применяется алгоритм Евклида: пара $(A,B)$ заменяется на
$(B,A\bmod B)$, пока второй элемент не станет нулём. Первый элемент
последней пары равен НОД.

\textbf{Библиотека \texttt{impl2}.}
Для решета выделяется массив из $B+1$ нулевых отметок. Для каждого
неотмеченного $p$ от $2$ до $\sqrt{B}$ отмечаются кратные $p$ начиная с $p^2$.
Меньшие кратные уже отмечены на шагах с меньшими делителями.
Затем подсчитываются неотмеченные числа от $\max(2,A)$ до $B$, и массив
освобождается. НОД ищется перебором от $\min(A,B)$ вниз. Первый делитель
обоих чисел и есть НОД.

\textbf{Подключение библиотек.}
CMake линкует \texttt{program1} с \texttt{impl1}, поэтому системный
загрузчик подключает её при запуске. В \texttt{program2} файл первой
библиотеки открывается из \texttt{main} по относительному пути.
Адреса функций ищутся по именам. Обе программы передают таблицу
указателей в общий цикл команд, который проверяет ввод и вызывает
выбранную функцию.

\textbf{Переключение.}
По команде \texttt{0} программа открывает другую библиотеку и ищет в ней
обе функции. Лишь после успешного поиска закрывает старую библиотеку
и заменяет её дескриптор и оба указателя. При ошибке открытия или поиска
прежняя реализация остаётся активной, если не возникла ошибка закрытия.
Активная библиотека закрывается при выходе.

CMake выбирает реализацию общего интерфейса \texttt{loader.h}.
В Linux открытие, поиск и закрытие выполняют \texttt{dlopen},
\texttt{dlsym} и \texttt{dlclose}. В Windows для тех же операций
используются \texttt{LoadLibraryW}, \texttt{GetProcAddress} и \texttt{FreeLibrary}.
